% Research notes: the stock engine's recurrent state under speculation (moved unchanged from the
% paper's Appendix D, except cross-references; owner: appendix editor).
\section{The stock engine under speculation}\label{notes:engine}
Appendix~\paperref{app:divergence} of the paper locates where stock configurations that ought to agree first produce different bits. This section holds how SGLang keeps the recurrent state under speculation and the tests of that path (H5); no claim of the paper depends on them.

\subsection{Recurrent state under speculation}\label{app:state}\label{sec:exactness-state}
Qwen's hybrid architecture makes state correctness part of verification. Reusable state includes attention KV, recurrent tensors, convolution history, positions and backend metadata. With state orientation $S\in\R^{d_v\times d_k}$, the gated-delta update can be written
\begin{equation}
S_t=\alpha_tS_{t-1}+u_tk_t^\top,\qquad u_t=\beta_t\bigl(v_t-\alpha_tS_{t-1}k_t\bigr),\label{eq:gdn}
\end{equation}
which is the Gated DeltaNet rule $S_t=S_{t-1}\alpha_t(I-\beta_tk_tk_t^\top)+\beta_tv_tk_t^\top$~\cite{gdn,deltanet} with the update vector made explicit. The update depends on the preceding state.
\begin{proposition}[Replay of valid derived records]\label{prop:gdn}
If $(\alpha_t,u_t,k_t)$ were recorded along a causal trajectory from $S_0$, applying the first $a$ records to that same $S_0$ reproduces $S_a$ in exact arithmetic:
\[S_a=\Bigl(\prod_{j=1}^{a}\alpha_j\Bigr)S_0+\sum_{t=1}^{a}\Bigl(\prod_{j=t+1}^{a}\alpha_j\Bigr)u_tk_t^\top.\]
\end{proposition}
\begin{proof}
Substitute the expansion at $a-1$ into~\eqref{eq:gdn}; the new gate multiplies every preceding term and the new outer product adds the last summand. Induction completes the proof.
\end{proof}
The qualification about $S_0$ matters: replaying a derived $u_t$ from a different state gives the wrong result even when the token IDs agree. The exact rational reference contains a one-dimensional witness, a record made at state zero and replayed from state two, which gives three where the genuine update gives one~\ev{inherited}. Replaying raw inputs with the true transition is a different operation. Regrouping the recurrence also does not preserve floating-point results (Appendix~\paperref{app:contracts} of the paper); SGLang's replay fold constrains its operation order and warp count so that it is a bitwise clone of the verify kernel's state update~\cite{gdncode}.

A cycle may emit a correction or bonus token before that token has been consumed into every cache, so the runtime must distinguish the \emph{emitted} prefix from the \emph{materialized} state prefix; treating both as one accepted length invites off-by-one state errors. Publication of speculative state should check that the request's slot and epoch are still current, and memory must not be reused while any stream or captured graph can still write it. An epoch check does not by itself prevent a stale kernel from writing reused memory; ownership and delayed reuse are separate obligations.

\paragraph{What SGLang does, and how it is tested (H5)} During verification SGLang computes the recurrent state and the convolution window after every draft position into per-position buffers and leaves the live state untouched; after acceptance a fused scatter copies the buffers of the last accepted position (for a tree, the last accepted node) into the request's slot (code reading at the pin). The emitted prefix includes the bonus token while the materialized state does not; the bonus token is the first input of the next cycle. Prefix reuse checkpoints the recurrent state every 256 tokens of sequence length into a ping-pong pair of slots, including checkpoints that fall inside a verify cycle, and a freshly allocated slot is cleared by the next forward on the same stream, after any in-flight verification by its previous owner. The MTP draft layer is a full-attention layer, and under greedy verification every committed token is the target's argmax, so an error in draft-side state could lower acceptance but not change the output.

Each point where a speculative cycle touches this state is tested against a reference that isolates it. Truncation by \code{max\_new\_tokens} and stop tokens at every index inside a verify cycle must reproduce the untruncated run's tokens on the same server. Checkpoints taken during speculative decoding, including in the cycle that finished the request, are restored for a continuation and compared with a cold prefill of the same prefix. Requests aborted mid-stream with a recurrent pool of exactly four slots force the next request onto the freed slot while the aborted request's last verification may still be in flight. Chunked prefill at 200 and 256 tokens is compared with an unchunked prefill. And the divergences of every MTP run from plain decoding are bucketed by the previous cycle's commit length, which fixes the draft position that cycle rejected. The first two tests are complete~\ev{targeted}, for MTP with three and five draft steps, top-$k$ 1, the prefix cache, overlap scheduler and CUDA graphs on, one request in flight and the first 40 prompts. Both tests compare output token IDs with the reference prefix; log-probabilities and state are not compared. Truncation leaves every truncated run token-identical to the untruncated run's prefix: 157 of 157 with three steps, keeping 1 to 4 tokens of the final cycle, and 240 of 240 with five steps, keeping 1 to 6. A stop token at every index of a verify cycle also leaves the output token-identical to the untruncated prefix, in 160 of 160 and 240 of 240 cases. In 120 and 200 of them the stop falls inside the draft block, so drafts after it were accepted and folded into the recurrent state before the stop was detected, and none of it changed an output token. Extending each stopped conversation with a new user turn, served warm from the restored checkpoint and cold after a flush, gives token-identical continuations in 139 of 160 and 210 of 240 cases; the rest diverge at exact ties (19 and 30) or within one BF16 step (2, both with three steps), consistent with the warm path's different prefill computation and with the history dependence above.

The same two tests for the tree, with three steps and top-$k$ 2, and for plain decoding also compare the top five log-probabilities~\ev{targeted}. Truncation leaves 159 of 159 tree runs and 39 of 39 plain runs identical in tokens and log-probabilities to the untruncated prefix, and a stop token at every index of a cycle leaves 160 of 160 and 40 of 40 identical up to the stop, 120 of the tree cases with the stop inside the draft block. Warm against cold continuation gives 142 of 160 identical for the tree (17 exact ties, 1 within one spacing) and 33 of 40 for plain decoding (7 exact ties); plain decoding has no speculative state, so such flips occur without it, consistent with the warm path's prefill and the prefix cache's history.

The bucketing by rejection position uses the pinned matrix at concurrency 1 against plain decoding~\ev{state-pinned}. At fragile positions, where plain decoding's top-two gap is at most 0.25 nats, the divergences per fragile position after each commit length of the previous cycle, full acceptance included, lie between 0.061 and 0.098 for MTP with one, three and five steps and the tree, and no length stands out. The rates are compatible with one common rate, with chi-square $p=0.75$, 0.72, 0.54 and 0.38; that is a failure to reject, not a proof of equal rates, and the test treats positions as independent although they cluster by prompt and a prompt's positions stop at its first divergence. No prompt-clustered analysis or bound on a per-length excess rate has been done. No divergence in any of the four MTP configurations falls at a non-fragile position, where a state error would show most plainly. The first verify cycle after the prefill follows no earlier cycle and is counted apart. Its rate is higher for five steps and the tree, 7 of 40 and 6 of 33 fragile positions against 0.077 and 0.081 per fragile position later, and a Fisher exact test against all later cycles gives $p=0.033$ and 0.049 (0.71 and 0.75 for one and three steps). That test is exploratory: it was chosen after looking at the data, it is not corrected for the four configurations and several cycle buckets examined, and the groups differ in more than state, since the first cycle follows the prefill chunk and every prompt contributes it, while later cycles count only prompts that have not yet diverged. It is a lead, not a finding. A test of it is declared before any of its data exist (\path{experiments/state_safety/README.md}, ``Declared follow-up''): on 960 fresh prompts disjoint from the 320, with frozen token IDs, it compares the first cycle with later ones for five steps and the tree against plain decoding at concurrency 1, and uses the same configurations at concurrency 1 against 32, which share the early exit of high-hazard prompts but do not compare verification with plain decoding, as a selection control. A first-cycle excess specific to speculation counts as supported only if the one-sided 95\% bootstrap lower bounds of both the pooled primary log odds ratio and its excess over the control's are above 0; otherwise the result is reported as inconclusive, not as evidence of no effect. Runs that fail the declared provenance, completeness or configuration checks are void, as is the analysis if every prompt is excluded, and a void result reports no statistic~\pend{state-firstcycle}. An effect that also differs between concurrency 1 and 32 would cancel in the control, and a supported result could be a benign difference in numerical path rather than a state error. A check of the recurrent state handed from the prefill to the first verify forward depends on its outcome. For three steps at batch 1 the cache-hashing tap already finds the caches entering the first verify forward, recurrent state included, identical to those entering plain decoding's first decode step, on 40 prompts (Appendix~\paperref{app:divergence} of the paper); five steps and the tree were not checked. Checkpoint reuse was tested with the prefix cache on~\ev{targeted}. Twelve thinking-mode prompts for MTP with three steps and 8 each for the tree and plain decoding generate 700 tokens, crossing the 256-token checkpoint interval during decode, and prefixes that end at a boundary or 1, 3 or 37 tokens past it are served for 48 tokens warm, after a flush and a regeneration of the original request, and cold, after a flush; tokens are compared, not log-probabilities. A finished request kept only its latest checkpoint, so a warm request restored a decode checkpoint only when its prefix extended past the last boundary the original crossed: in 9 of 100 cases for MTP, 6 of 68 for the tree and 6 of 68 for plain decoding. Where it did, warm and cold gave the same tokens in 9 of 9, 6 of 6 and 4 of 6 cases, the other two at exact ties. A checkpoint taken in the verify cycle that finished the request was restored in every case, and warm and cold agreed in 45 of 53 cases for MTP and 19 of 21 for the tree, the rest at exact ties. Warm and cold reach the boundary state by different computations, decode or verify forwards against a chunked prefill, so bitwise equality is not expected; every difference is an exact tie or within one BF16 step, and plain decoding, which has no speculative state, shows them too, in 9 of 68 cases against 21 of 153 for MTP and 8 of 89 for the tree. No difference points to a wrong restored state, but few decode checkpoints were restored. Aborts with slot reuse, chunked prefill and run-to-run repeats are \pend{state-targeted}.
